\section{Ejemplo de cálculo de calor para un cilindro}\label{Ap:Temp_cilindro}

Si T1 > T2 el flujo de calor radial se da desde el interior hasta el exterior del contenedor.
\begin{center}
    
    $Q = -AK\frac{\Delta T}{\Delta \varepsilon} $
\end{center}

Se tienen dos áreas de las paredes de un cilindro hueco:
\begin{center}
    
$A_{\text{int}} = 2\pi r_1 L$ 
\end{center}
\begin{center}
    
$A_{\text{ext}} = 2\pi r_2 L$
\end{center}
Se toma un radio de referencia r que varía desde r1 hasta r2 

Sustituyendo el área de la pared de un cilindro, tomando en cuenta el radio de referencia r.


\[
Q = -(2\pi rL) \frac{dT}{dr}
\]

\[
\frac{-Q}{2\pi L} \int_{r_1}^{r_2} \frac{dr}{r} = k \int_{T_1}^{T_2} dT
\]

\[
\frac{-Q}{2\pi L} (\ln r_2 - \ln r_1) = k(T_2 - T_1)
\]

\[
Q = \frac{k2\pi L(T_1 - T_2)}{\ln(r_2/r_1)}
\]

Para las tapas del cilindro:

\[
Q = \frac{k2\pi r^2 (T_2 - T_1)}{\varepsilon}
\]